\section{Modelos avanzados}

\subsection{Arquitecturas evaluadas}

Se conservaron cuatro backbones con pesos ImageNet. EfficientNet-B0 escala
profundidad, ancho y resolución de forma conjunta \citep{tan2019efficientnet};
ShuffleNetV2-x1.0 prioriza operaciones eficientes y menor movimiento de memoria
\citep{ma2018shufflenet}; MobileNetV3 Small fue diseñado mediante búsqueda para
dispositivos con recursos limitados \citep{howard2019mobilenetv3}; FastViT-T8
combina bloques convolucionales e ideas de transformers visuales con
reparametrización estructural \citep{vasu2023fastvit}.

Los cuatro recibieron exactamente las mismas filas de entrenamiento y
validación y la configuración de cabeza seleccionada con EfficientNet-B0.
Transferir esa misma cabeza metodológica evita darle a cada arquitectura una
búsqueda de 30 trials diferente, aunque también favorece parámetros que fueron
escogidos sobre B0. Los resultados miden la calidad de la representación
congelada, no un ranking definitivo después de fine-tuning.

El registro de código también admite EfficientNet-Lite0, EfficientNet-B4,
MobileNetV3 Large y GhostNetV2-100. Registrado significa que el constructor
puede instanciarlo; evaluado significa que produjo embeddings para las 33\,433
imágenes y un CSV de predicciones en este protocolo. Solo los cuatro modelos
de la Tabla~\ref{tab:model-comparison} cumplen la segunda condición.

\subsection{Comparación experimental}

\begin{table}[H]
  \centering
  \caption{Backbones evaluados con la misma cabeza optimizada. MB representa el
  tamaño FP32 estimado de parámetros del backbone; ms es la mediana de 30
  inferencias de una imagen en CPU.}
  \label{tab:model-comparison}
  \resizebox{\textwidth}{!}{\input{tables/model_comparison}}
\end{table}

\begin{figure}[H]
  \centering
  \includegraphics[width=0.86\textwidth]{model_comparison.pdf}
  \caption{Relación entre macro-F1, tamaño FP32 y latencia del backbone. El
  color de cada punto corresponde a la latencia CPU.}
  \label{fig:model-comparison}
\end{figure}

La latencia separó dos grupos aun antes de mirar la calidad. MobileNetV3 Small
registró 5.2 ms y ShuffleNetV2-x1.0 9.5 ms por backbone; FastViT-T8 y
EfficientNet-B0 rondaron 23.5 y 24.3 ms. Estas mediciones excluyen lectura de
archivo, preprocesamiento y cabeza, y proceden de un Intel Core i7-12800HX. No
son una predicción de tiempo Android.

EfficientNet-B0 conservó el primer lugar: 0.9382 de accuracy y 0.8887 de
macro-F1. ShuffleNet quedó segundo con 0.9264 y 0.8694. La diferencia de
macro-F1 fue 0.0193, mientras ShuffleNet utilizó cerca del 31\% del tamaño FP32
de parámetros de B0 y el 39\% de su latencia. Se repite así el trade-off de la
Etapa 1: B0 ofrece la mejor calidad; ShuffleNet pierde poco y reduce de forma
marcada el costo.

MobileNetV3 Small fue todavía más rápido, 5.2 ms, pero su macro-F1 bajó a
0.8651. Frente a ShuffleNet ahorra aproximadamente 4.2 ms en esta CPU a cambio
de 0.0043 de macro-F1. Ese intercambio puede ser razonable en un teléfono
lento, aunque requiere una medición Android antes de elegirlo. FastViT-T8 no
justificó aquí su costo: obtuvo el macro-F1 menor, 0.8593, con latencia similar
a EfficientNet-B0 y más parámetros que los otros dos modelos móviles.

La comparación tampoco respalda la idea de que una arquitectura más reciente
deba ganar automáticamente. El backbone híbrido fue competitivo en accuracy
(0.9194), pero su promedio por clase dejó ver que la representación no separó
las minorías tan bien como B0. Por eso EfficientNet-B0 pasó a la fase de
ensemble como mejor candidato individual y no por continuidad histórica.
